% ============================================================================
\bigskip
\section{Phương pháp đề xuất và thiết kế hệ thống}
\label{sec:method}
% ============================================================================

\subsection{Tổng quan kiến trúc}

Hệ thống được tổ chức thành bốn tầng (Hình~\ref{fig:architecture}). \textbf{Tầng dữ
liệu} nạp UIT-ViQuAD~2.0, kiểm tra rò rỉ và khả năng chấm điểm, rồi gán nhãn phục vụ
phân tích. \textbf{Tầng mô hình} gồm bốn bộ dự đoán cùng tuân theo một giao diện
\code{Predictor.predict(context, question) -> str}, nên harness đánh giá xử lý baseline
và Transformer hoàn toàn như nhau. \textbf{Tầng đánh giá} chạy bộ dự đoán, tính metric,
tách nhóm và ghi kết quả kèm thông tin truy vết. \textbf{Tầng sản phẩm} sinh hình vẽ từ
tệp kết quả, kiểm tra giả thuyết và phục vụ ứng dụng web.

\begin{figure}[htbp]
\centering
\resizebox{\textwidth}{!}{%
\begin{tikzpicture}[
  font=\small,
  box/.style={draw=black!45,rounded corners=3pt,fill=white,align=center,minimum height=1.0cm,inner sep=4pt},
  layer/.style={draw=none,rounded corners=6pt,inner sep=8pt},
  lbl/.style={font=\footnotesize\bfseries,anchor=south west,inner sep=2pt,fill=white,rounded corners=2pt},
  arr/.style={-{Latex[length=2.2mm]},black!65,thick},
  gate/.style={draw=badred,thick,fill=red!5,rounded corners=2pt,font=\footnotesize,align=center,inner sep=4pt,minimum height=1.0cm}]

  % Tầng dữ liệu
  \node[box] (hf) at (0,0) {HuggingFace\\\code{taidng/UIT-ViQuAD2.0}};
  \node[box,right=0.8cm of hf] (data) {\code{data.py}\\parse SQuAD~2.0 $\to$ \code{Example}};
  \node[gate,right=0.8cm of data] (gates) {\code{assert\_no\_leakage}\\\code{assert\_gradeable}};
  \node[box,right=0.8cm of gates] (tag) {\code{tagging.py}\\độ dài, loại câu hỏi};
  \draw[arr] (hf) -- (data); \draw[arr] (data) -- (gates); \draw[arr] (gates) -- (tag);

  % Tầng mô hình
  \node[box] (tfidf) at (0,-2.6) {\code{baseline\_tfidf.py}\\TF-IDF + cosine};
  \node[box,right=0.5cm of tfidf] (xlmr) {XLM-R squad2\\(zero-shot)};
  \node[box,right=0.5cm of xlmr] (mbert) {mBERT + QA\\(fine-tune)};
  \node[box,right=0.5cm of mbert] (viso) {ViSoBERT + QA\\(fine-tune)};
  \node[box,right=0.8cm of viso,fill=orange!10] (ft) {\code{finetune.py}\\\code{training.py}, \code{features.py}};
  \draw[arr] (ft.west) -- (viso.east);
  \node[box,fill=uitblue!8,minimum width=9.2cm] (tqa) at ($(mbert.south)+(-0.2,-1.0)$)
     {\code{transformer\_qa.py} + \code{windowing.py}: cửa sổ trượt, offset, chọn span, ngưỡng null};
  \foreach \m in {xlmr,mbert,viso} { \draw[black!45,densely dashed] (\m.south) -- (\m.south |- tqa.north); }

  % Tầng đánh giá
  \node[box,fill=okgreen!10,minimum width=7.6cm] (eval) at ($(tqa.south)+(-1.2,-1.6)$)
     {\code{evaluate.py} + \code{metrics.py}\\EM, F1, answerable/impossible, breakdown, provenance};
  \node[box,right=0.9cm of eval] (json) {\code{results/eval\_*.json}\\commit · thời điểm · thiết bị · $n$};
  \draw[arr] (eval) -- (json);

  % Tầng sản phẩm
  \node[box] (fig) at ($(eval.south)+(-2.6,-1.7)$) {\code{make\_figures.py}\\hình vẽ báo cáo};
  \node[box,right=0.6cm of fig] (hyp) {\code{check\_hypotheses.py}\\đối chiếu giả thuyết};
  \node[box,right=0.6cm of hyp] (app) {\code{streamlit\_app.py}\\demo web};
  \coordinate (jbus) at ($(json.south)+(0,-0.55)$);
  \draw[arr] (json.south) -- (jbus) -| (hyp.north);
  \draw[arr] (jbus) -| (fig.north);
  \draw[arr] (ft.south) |- ($(app.north)+(0,0.8)$) -- (app.north);

  % Luồng chính
  \coordinate (dbus) at ($(tag.south)+(0,-0.7)$);
  \draw[arr] (tag.south) -- (dbus) -| (ft.north);
  \draw[arr] (dbus) -| (tfidf.north);
  \draw[arr] (tfidf.south) |- (eval.west);
  \draw[arr] (tqa.south) -- (tqa.south |- eval.north);

  % Nền các tầng
  \begin{scope}[on background layer]
    \node[layer,fill=uitblue!4,fit=(hf)(tag)] (L1) {};
    \node[layer,fill=orange!4,fit=(tfidf)(ft)(tqa)] (L2) {};
    \node[layer,fill=okgreen!4,fit=(eval)(json)] (L3) {};
    \node[layer,fill=black!4,fit=(fig)(app)] (L4) {};
  \end{scope}
  \node[lbl,text=uitblue]    at (L1.north west) {TẦNG DỮ LIỆU};
  \node[lbl,text=warnorange] at (L2.north west) {TẦNG MÔ HÌNH --- giao diện chung \code{Predictor.predict(context, question)}};
  \node[lbl,text=okgreen]    at (L3.north west) {TẦNG ĐÁNH GIÁ};
  \node[lbl,text=black!70]   at (L4.north west) {TẦNG SẢN PHẨM};
\end{tikzpicture}}
\caption{Kiến trúc tổng thể của hệ thống. Ô viền đỏ là cổng kiểm tra làm dừng chương trình khi vi phạm}
\label{fig:architecture}
\end{figure}

\subsection{Nguyên tắc thiết kế: bốn bất biến}
\label{sec:invariants}

Thiết kế của dự án là phản ứng trực tiếp với thất bại của phiên bản trước đó của chính
đồ án này: một bản kiểm toán nội bộ phát hiện báo cáo cũ trình bày những con số chưa từng
được sinh ra bởi bất kỳ lần chạy nào, và các con số mâu thuẫn giữa README, báo cáo và
ứng dụng. Kết luận rút ra là đây không phải lỗi cẩu thả mà là \emph{lỗi kiến trúc}: không
có thước đo đáng tin nào được xác lập trước khi có số để báo cáo. Vì vậy dự án mới đặt
bốn bất biến; vi phạm bất kỳ bất biến nào đều làm test thất bại hoặc làm dừng chương
trình (Bảng~\ref{tab:invariants}).

\begin{table}[htbp]
\centering
\caption{Bốn bất biến của hệ thống và cơ chế thực thi}
\label{tab:invariants}
\renewcommand{\arraystretch}{1.2}
\small
\begin{tabularx}{\textwidth}{@{}c l L L@{}}
\toprule
\textbf{\#} & \textbf{Bất biến} & \textbf{Phát biểu} & \textbf{Cơ chế thực thi} \\
\midrule
1 & Trích xuất & Mọi \code{predict()} trả về chuỗi con của \code{context} hoặc chuỗi rỗng & Test thuộc tính cho các bộ dự đoán; cắt chuỗi gốc theo offset, không dùng \code{decode()} \\
2 & Không rò rỉ & Không đoạn văn nào xuất hiện ở cả hai split & \code{assert\_no\_leakage()} chạy trước khi huấn luyện, \emph{raise} khi vi phạm \\
3 & Truy vết được & Mọi con số gắn với \code{commit}, \code{timestamp}, \code{device}, \code{split}, $n$ & Harness ghi các trường này vào mọi tệp kết quả \\
4 & Có $n$ kèm số & Không bảng nào có số thiếu cỡ mẫu; nhóm $n < 30$ bị gắn cờ & Trường \code{count} và \code{unreliable} trong mọi breakdown \\
\bottomrule
\end{tabularx}
\end{table}

\subsection{Tầng dữ liệu}
\label{sec:datalayer}

\subsubsection{Phân biệt ba loại câu hỏi}

Mỗi câu hỏi được biểu diễn bằng một đối tượng \code{Example} bất biến (frozen dataclass)
gồm \code{qid}, \code{question}, \code{context}, \code{title}, \code{answers},
\code{answer\_start} và \code{is\_impossible}. Điểm thiết kế quan trọng nhất là thuộc tính
\code{is\_gradeable}, phân biệt ba trường hợp mà một cách cài đặt ngây thơ sẽ gộp làm một
(Bảng~\ref{tab:gradeable}).

\begin{table}[htbp]
\centering
\caption{Ba trường hợp nhãn và khả năng chấm điểm}
\label{tab:gradeable}
\renewcommand{\arraystretch}{1.2}
\small
\begin{tabularx}{\textwidth}{@{}l c c L@{}}
\toprule
\textbf{Tình huống} & \textbf{Gold} & \textbf{Chấm được?} & \textbf{Đáp án đúng} \\
\midrule
Câu có đáp án & không rỗng & có & Chuỗi gold \\
\code{is\_impossible = True} & rỗng & có & Chuỗi rỗng \\
Gold rỗng, \code{is\_impossible = False} & rỗng & \textbf{không} & Không biết --- nhãn đã bị lược bỏ (blind split) \\
\bottomrule
\end{tabularx}
\end{table}

Trường hợp thứ ba chính là toàn bộ test split của UIT-ViQuAD~2.0 (Mục~\ref{sec:blind}).
Nếu suy ra \code{is\_impossible} từ việc gold rỗng, cả 7.301 câu test sẽ bị coi là câu
không có đáp án, và một mô hình \emph{luôn trả về chuỗi rỗng} sẽ đạt EM 100\% --- một con
số trông hoàn toàn hợp lệ trong báo cáo. Hàm \code{assert\_gradeable()} được gọi ở đầu
harness đánh giá và từ chối chấm một split như vậy.

\subsubsection{Chống rò rỉ theo đoạn văn}

Mỗi đoạn văn trong UIT-ViQuAD có nhiều câu hỏi. Nếu chia dữ liệu theo \emph{câu hỏi}, mô
hình có thể được huấn luyện trên một số câu hỏi của một đoạn văn rồi được đánh giá trên
các câu hỏi khác của chính đoạn văn đó --- rò rỉ theo nhóm (group leakage), cho kết quả
cao giả tạo. Đơn vị chia và kiểm tra rò rỉ trong dự án là \textbf{đoạn văn}
(\code{context}), không phải bài viết (\code{title}): tập train có 138 bài viết nhưng
4.101 đoạn văn, và nhầm hai đơn vị này làm vô hiệu hoá chính cơ chế bảo vệ. Hàm
\code{split\_by\_context()} cài đặt chia nhóm có seed (tương đương \code{GroupShuffleSplit}
với nhóm là đoạn văn) cho trường hợp cần tự chia; với train/validation có sẵn,
\code{assert\_no\_leakage()} kiểm tra giao của hai tập đoạn văn phải rỗng và được gọi trước
khi bắt đầu mọi lần fine-tune.

\subsubsection{Lấy mẫu con tái lập được và không thiên lệch}
\label{sec:subset}

Để giữ thời gian đánh giá trong tầm kiểm soát, các thực nghiệm dùng mẫu con 500 câu. Hàm
\code{reproducible\_subset(examples, n, seed=42)} lấy mẫu \emph{ngẫu nhiên có seed} thay vì
lấy $n$ câu đầu tệp. Lựa chọn này được đưa ra sau khi phát hiện một lỗi thực tế: câu hỏi
trong tệp được sắp theo bài viết, nên $n$ câu đầu chỉ thuộc vài chủ đề. Cùng một checkpoint
mBERT (epoch 2) cho \textbf{EM \curveMbertFirstEMII{}} trên 300 câu đầu tệp nhưng
\textbf{EM \curveMbertValEMII{}} trên 300 câu ngẫu nhiên --- chênh \curveMbertFirstGap{}
điểm chỉ do cách lấy mẫu. Hàm này được dùng chung cho cả
đường cong huấn luyện lẫn bảng kết quả cuối, nên hai nguồn số liệu so sánh được với nhau.

\subsection{Chuẩn hoá và metric cho tiếng Việt}
\label{sec:metric-design}

Metric được cài đặt \emph{trước} mọi mô hình, theo nguyên tắc ``thước đo phải được kiểm
chứng trước vật được đo'', và đạt độ phủ kiểm thử 100\%. Bộ đánh giá SQuAD gốc được điều
chỉnh bằng ba quyết định, mỗi quyết định được ``ghim'' bằng test để người sau không vô
tình ``sửa cho giống SQuAD'' (Bảng~\ref{tab:metric-decisions}).

\begin{table}[htbp]
\centering
\caption{Ba quyết định chuẩn hoá đặc thù tiếng Việt}
\label{tab:metric-decisions}
\renewcommand{\arraystretch}{1.22}
\small
\begin{tabularx}{\textwidth}{@{}c l L L@{}}
\toprule
& \textbf{Quyết định} & \textbf{Lý do} & \textbf{Hành vi được test ghim} \\
\midrule
A & Không loại bỏ mạo từ & Tiếng Việt không có mạo từ; ``các'', ``những'', ``con'' là loại từ và có thể thuộc đáp án & ``các tỉnh'' giữ nguyên ``các'' sau chuẩn hoá \\
B & Tách token theo khoảng trắng & Khớp cách đánh giá thông dụng của ViQuAD; không phụ thuộc và không nhiễm lỗi bộ tách từ & ``Hà Nội'' là 2 token \\
C & Giữ dấu tiếng Việt & ``hoà'' $\ne$ ``hoa''; bỏ dấu làm sai EM hàng loạt & EM(``hoà'', ``hoa'') $= 0$ \\
\bottomrule
\end{tabularx}
\end{table}

Hai chi tiết cài đặt khác cũng được test: dấu câu được nhận diện theo nhóm Unicode
\code{P*} thay vì \code{string.punctuation}, để bắt được các ký tự như ``–'', ``“'', ``…''
thường gặp trong Wikipedia; và F1 dùng \code{Counter} (đa tập) thay vì \code{set}, để token
lặp được đếm đúng. Hàm tổng hợp \code{evaluate()} \emph{raise} \code{KeyError} nếu thiếu dự
đoán cho bất kỳ câu nào: âm thầm bỏ qua câu thiếu sẽ khiến EM được tính trên tập con trong
khi báo cáo ghi $n$ đầy đủ.

\subsection{Baseline TF-IDF}
\label{sec:baseline}

Thuật toán~\ref{alg:tfidf} mô tả baseline. Hai chi tiết đáng chú ý. Thứ nhất, bộ tách
câu dùng biểu thức chính quy \verb/(?<!\d)[.!?]+(?:\s+|$)/: dấu chấm đứng sau chữ số
không kết thúc câu, tránh cắt ``8.5 triệu'' thành hai câu; dấu kết thúc câu được giữ lại
để câu trả về vẫn là chuỗi con của đoạn văn (bất biến~1). Thứ hai, TF-IDF được
\code{fit} trên chính các câu của đoạn văn cùng câu hỏi --- truy hồi trong tài liệu, không
dùng kho văn bản ngoài, nên không có rò rỉ. N-gram $(1,2)$ giúp bắt các từ hai âm tiết như
``thủ đô'', ``Hà Nội''.

\begin{algorithm}[H]
\caption{Baseline truy hồi câu bằng TF-IDF}
\label{alg:tfidf}
\KwIn{đoạn văn $C$, câu hỏi $Q$}
\KwOut{một câu của $C$ (chuỗi con của $C$)}
$\mathcal{S} \leftarrow \textsc{TáchCâu}(C)$\;
\If{$|\mathcal{S}| = 0$}{\Return chuỗi rỗng\;}
\If{$|\mathcal{S}| = 1$}{\Return $\mathcal{S}_1$\;}
$M \leftarrow \textsc{TfidfVectorizer}(\text{ngram}=(1,2)).\textsc{fit\_transform}(\mathcal{S} \cup \{Q\})$\;
$\mathbf{s} \leftarrow \cos(M_Q, M_{\mathcal{S}})$\tcp*{độ tương đồng với từng câu}
\Return $\mathcal{S}_{\arg\max \mathbf{s}}$\;
\end{algorithm}

\subsection{Suy luận với mô hình Transformer}
\label{sec:inference}

\subsubsection{Quy trình}

Lớp \code{TransformerQA} bọc \code{AutoModelForQuestionAnswering} thành một
\code{Predictor}. Khi khởi tạo, lớp này \emph{từ chối} mọi mô hình không có fast tokenizer
thay vì chạy rồi cho kết quả sai. Thuật toán~\ref{alg:infer} tóm tắt quy trình suy luận.

\begin{algorithm}[H]
\caption{Suy luận hỏi đáp trích xuất với cửa sổ trượt}
\label{alg:infer}
\KwIn{đoạn văn $C$, câu hỏi $Q$, mô hình $f$, \code{max\_length} $=384$, \code{doc\_stride} $=128$, $L_{\max}$, $\tau$}
\KwOut{chuỗi con của $C$, hoặc chuỗi rỗng}
$\mathcal{W} \leftarrow \textsc{MakeWindows}(Q, C)$\tcp*{offset tuyệt đối trong $C$}
$\text{best} \leftarrow -\infty$;\quad $\text{span} \leftarrow \varnothing$\;
\ForEach{cửa sổ $w \in \mathcal{W}$}{
  $(S, E) \leftarrow f(w)$\;
  $(i, j) \leftarrow$ span tốt nhất theo công thức~\eqref{eq:span}\;
  \If{$S_i + E_j \le S_{\mathrm{CLS}} + E_{\mathrm{CLS}} + \tau$}{\Continue\tcp*{cửa sổ nói ``không có đáp án''}}
  \If{$S_i + E_j > \text{best}$}{$\text{best} \leftarrow S_i + E_j$;\quad $\text{span} \leftarrow (w.\mathrm{start}_i,\, w.\mathrm{end}_j)$\;}
}
\If{$\text{span} = \varnothing$}{\Return chuỗi rỗng\;}
\Return $C[\text{span}]$\tcp*{cắt chuỗi gốc, không dùng \code{decode()}}
\end{algorithm}

\subsubsection{Tự cài đặt cửa sổ trượt}
\label{sec:window-bug}

Cách thông thường để tạo cửa sổ là tham số \code{return\_overflowing\_tokens=True} của
tokenizer. Trong quá trình kiểm thử, nhóm phát hiện rằng với phiên bản
\code{transformers 5.17.0} được dùng, cơ chế này \textbf{chỉ sinh 2 cửa sổ} bất kể độ dài
đoạn văn: đoạn văn 210, 420, 700 và 1.400 token đều sinh 2 cửa sổ, trong khi theo công
thức~\eqref{eq:nwin} phải là 2, 5, 8 và 16. Phần đuôi đoạn văn bị cắt \emph{âm thầm} ---
không ngoại lệ, không cảnh báo --- nên đáp án nằm cuối đoạn văn dài sẽ không bao giờ được
tìm thấy khi suy luận, và không bao giờ được gán nhãn khi huấn luyện. Lỗi được phát hiện bởi
test \code{test\_at\_least\_one\_window\_of\_long\_context\_finds\_the\_answer}.

Dự án tự cài hàm \code{make\_windows()} (Hình~\ref{fig:windows}): tokenize đoạn văn một
lần để lấy offset của từng token trong chuỗi gốc, cắt theo ngân sách token $B$ với bước
$B - \texttt{doc\_stride}$, tokenize từng cặp (câu hỏi, đoạn cắt) với
\code{truncation="only\_second"} để không bao giờ cắt câu hỏi, rồi \emph{dịch offset của
đoạn cắt về hệ toạ độ của đoạn văn gốc}. Ở \code{max\_length} $=384$, ảnh hưởng thực tế trên
tập validation là nhỏ (2 trên 557 đoạn văn cần nhiều hơn 2 cửa sổ), nhưng vì đây là lỗi
đúng--sai âm thầm, nó được sửa tận gốc thay vì chấp nhận.

\begin{figure}[htbp]
\centering
\resizebox{\textwidth}{!}{%
\begin{tikzpicture}[font=\small,
  win/.style={draw,rounded corners=2pt,minimum height=0.55cm,inner sep=2pt},
  arr/.style={-{Latex[length=2mm]},black!60}]
  \fill[black!6] (0,0) rectangle (14,0.6);
  \draw[black!40] (0,0) rectangle (14,0.6);
  \node at (7,0.3) {Đoạn văn gốc: $N$ token, mỗi token mang offset ký tự tuyệt đối};
  \draw[decorate,decoration={brace,amplitude=4pt}] (0,0.75) -- (3.6,0.75) node[midway,above=4pt,font=\footnotesize] {bước $= B - \texttt{doc\_stride}$};
  \node[win,draw=uitblue,fill=uitblue!10,minimum width=5.6cm,anchor=west] (w1) at (0,-0.75) {cửa sổ 1: \code{[CLS] Q [SEP]} đoạn $[0, B)$ \code{[SEP]}};
  \node[win,draw=okgreen,fill=okgreen!10,minimum width=5.6cm,anchor=west] (w2) at (3.6,-1.5) {cửa sổ 2};
  \node[win,draw=warnorange,fill=orange!10,minimum width=5.6cm,anchor=west] (w3) at (7.2,-2.25) {cửa sổ 3};
  \node[win,draw=badred,fill=red!8,minimum width=3.2cm,anchor=west] (w4) at (10.8,-3.0) {cửa sổ 4 (phần đuôi)};
  \draw[decorate,decoration={brace,amplitude=4pt,mirror}] (3.6,-1.85) -- (5.6,-1.85) node[midway,below=4pt,font=\footnotesize] {chồng lấp $=$ \code{doc\_stride}};
  \node[align=left,font=\footnotesize,text=badred,anchor=west] at (0,-3.1)
     {\textbf{\code{return\_overflowing\_tokens}} (transformers 5.17.0):\\ dừng sau cửa sổ 2 $\Rightarrow$ cửa sổ 3, 4 bị bỏ \emph{âm thầm}};
\end{tikzpicture}}
\caption{Cửa sổ trượt tự cài đặt. Offset của mỗi cửa sổ được dịch về hệ toạ độ của đoạn văn gốc để cắt đáp án trực tiếp từ chuỗi gốc}
\label{fig:windows}
\end{figure}

\subsubsection{Ba lỗi chọn span được chặn bằng test}

Hàm thuần \code{select\_best\_span()} (công thức~\eqref{eq:span}--\eqref{eq:null}) được kiểm
thử bằng logit tổng hợp, không cần tải mô hình. Các test chặn ba lỗi kinh điển không gây
ngoại lệ: (i) span lấy từ vùng câu hỏi, (ii) span đảo ngược khi $\arg\max$ độc lập cho
điểm đầu và điểm cuối, (iii) lấy lại chuỗi bằng \code{decode()} làm mất dấu. Hàm
\code{decode\_span()} \emph{raise} khi offset không hợp lệ thay vì trả chuỗi rỗng, vì một
span sai là lỗi cần sửa, không phải ``không tìm thấy đáp án''.

\subsection{Fine-tuning}
\label{sec:finetune}

\subsubsection{Chuyển vị trí đáp án từ ký tự sang token}

Hàm \code{prepare\_train\_features()} tạo các cửa sổ bằng cùng \code{make\_windows()} dùng khi
suy luận, rồi với mỗi cửa sổ tìm cặp token bao trọn khoảng ký tự
$[\text{answer\_start},\ \text{answer\_start} + |\text{answer}|)$. Nếu đáp án không nằm trọn
trong cửa sổ, hoặc câu hỏi không có đáp án, nhãn được gán về \code{[CLS]}. Đây là chỗ sai
âm thầm nguy hiểm nhất của toàn pipeline: nếu ánh xạ lệch, mô hình học nhãn sai mà
\emph{loss vẫn giảm bình thường}. Vì vậy bất biến được kiểm thử không phải ``chạy không
lỗi'', mà là \emph{giải mã ngược nhãn token phải ra đúng chuỗi đáp án vàng}.

\subsubsection{Siêu tham số}

Bảng~\ref{tab:hparams} liệt kê cấu hình fine-tune, lấy trực tiếp từ trường \code{config}
trong tệp đường cong huấn luyện và nhật ký huấn luyện. Tốc độ học $3\cdot10^{-5}$ và 2 epoch
là cấu hình thông dụng khi fine-tune mô hình họ BERT cho QA. ViSoBERT dùng tốc độ học, số
epoch và $L_{\max}$ lớn hơn, sau quá trình chẩn đoán được trình bày ở
Mục~\ref{sec:visobert-analysis}. Không có tìm kiếm siêu tham số có hệ thống.

\begin{table}[htbp]
\centering
\caption{Cấu hình fine-tune (nguồn: \code{results/training\_curve\_*.json}, \code{results/finetune\_*.log})}
\label{tab:hparams}
\renewcommand{\arraystretch}{1.15}
\footnotesize
\setlength{\tabcolsep}{5pt}
\begin{tabular}{@{}l r r l@{}}
\toprule
\textbf{Tham số} & \textbf{mBERT} & \textbf{ViSoBERT} & \textbf{Ghi chú} \\
\midrule
Dữ liệu huấn luyện & 28.454 câu & 28.454 câu & toàn bộ train \\
Số feature (cửa sổ) & 30.540 & 40.984 & ViSoBERT phân mảnh nhiều hơn \\
Số batch mỗi epoch & 2.545 & 3.416 & batch size 12 \\
Số epoch & 2 & 3 & \\
Tích luỹ gradient & 2 & 2 & batch hiệu dụng 24 \\
Tốc độ học $\eta_{\max}$ & $3\cdot10^{-5}$ & $5\cdot10^{-5}$ & \\
Warmup / weight decay & 0,1 / 0,01 & 0,1 / 0,01 & warmup theo tỉ lệ tổng số bước \\
\code{max\_length} / \code{doc\_stride} & 384 / 128 & 384 / 128 & \\
$L_{\max}$ (độ dài đáp án tối đa) & 30 & 64 & phụ thuộc tokenizer \\
Cắt gradient & 1,0 & 1,0 & chuẩn L2 \\
Đánh giá mỗi epoch & 300 câu & 300 câu & ngẫu nhiên, seed 42 \\
Seed & 42 & 42 & \\
Độ chính xác số & fp32 & fp32 & không dùng fp16 kiểu CUDA trên MPS \\
\bottomrule
\end{tabular}
\end{table}

\subsubsection{Chọn epoch và chẩn đoán đường cong}

Sau mỗi epoch, checkpoint được lưu và đánh giá trên 300 câu validation ngẫu nhiên (seed 42)
bằng \emph{chính} pipeline suy luận dùng cho kết quả cuối. Epoch có F1 validation cao nhất
được chọn làm mô hình cuối (hoà thì chọn epoch sớm hơn) --- chọn theo F1 chứ không theo
loss, vì loss là đại lượng trên tập train. Toàn bộ logic ra quyết định nằm trong module
\code{training.py} dưới dạng hàm thuần, được test trong vài mili-giây thay vì phải chạy một
epoch 35 phút mới biết đúng sai:

\begin{itemize}
  \item \code{compute\_schedule()} tính số bước cập nhật theo đơn vị \emph{lần cập nhật}
        (sau tích luỹ gradient), không phải theo batch.
  \item \code{detect\_overfitting()} trả về ba tín hiệu: \emph{overfitting} (F1 validation
        đạt đỉnh rồi giảm trong khi loss train vẫn giảm), \emph{still\_improving} (F1 vẫn
        tăng ở epoch cuối --- thiếu epoch) và \emph{degenerate\_collapse}.
  \item \emph{Degenerate collapse} được phát hiện khi $|\mathrm{EM} - \mathrm{F1}| \le 0{,}5$
        ở mọi epoch. F1 cho điểm bán phần nên bình thường cao hơn EM; hai giá trị trùng nhau
        nghĩa là mọi câu chỉ đúng hoàn toàn hoặc sai hoàn toàn --- dấu hiệu mô hình luôn trả
        về chuỗi rỗng (Mục~\ref{sec:metrics-theory}).
\end{itemize}

\subsection{Lựa chọn mô hình như một thiết kế thí nghiệm}
\label{sec:models}

Bốn mô hình không được chọn để ``thử nhiều cho chắc'', mà để mỗi bước trong chuỗi trả lời
một câu hỏi khoa học riêng (Bảng~\ref{tab:model-design}).

\begin{table}[htbp]
\centering
\caption{Bốn mô hình và câu hỏi khoa học tương ứng}
\label{tab:model-design}
\renewcommand{\arraystretch}{1.22}
\small
\begin{tabularx}{\textwidth}{@{}l l c L@{}}
\toprule
\textbf{Mô hình} & \textbf{Checkpoint} & \textbf{Fine-tune} & \textbf{Câu hỏi khoa học} \\
\midrule
TF-IDF & --- & không & Bài toán có giải được bằng so khớp từ khoá thuần tuý? \\
XLM-R & \code{deepset/xlm-roberta-base-squad2} & không & Chuyển giao từ SQuAD~2.0 tiếng Anh sang tiếng Việt được bao nhiêu? \\
mBERT & \code{bert-base-multilingual-cased} & có & Fine-tune trên dữ liệu tiếng Việt trong miền thêm được bao nhiêu? \\
ViSoBERT & \code{uitnlp/visobert} & có & Tiền huấn luyện chuyên tiếng Việt đóng góp bao nhiêu? \\
\bottomrule
\end{tabularx}
\end{table}

Cần lưu ý rằng chuỗi này \emph{không} phải ablation thuần tuý: XLM-R và mBERT khác nhau
đồng thời ở mô hình nền và ở dữ liệu fine-tune; mBERT và ViSoBERT khác nhau đồng thời ở
ngôn ngữ, miền tiền huấn luyện, cỡ từ vựng và siêu tham số huấn luyện. Các kết luận ở
Mục~\ref{sec:experiments} vì vậy được phát biểu ở mức ``cấu hình A so với cấu hình B'',
không quy toàn bộ chênh lệch cho một yếu tố duy nhất.

\subsubsection{Vì sao không dùng PhoBERT}
\label{sec:phobert}

Đề tài gợi ý ``PhoBERT encoder + QA head'', và kế hoạch ban đầu của nhóm dùng
\code{vinai/phobert-base-v2}. Kiểm tra trực tiếp cho thấy điều này không khả thi với
pipeline trích xuất dựa trên offset (Bảng~\ref{tab:phobert}).

\begin{table}[htbp]
\centering
\caption{Kiểm tra các encoder tiếng Việt ứng viên}
\label{tab:phobert}
\renewcommand{\arraystretch}{1.18}
\small
\begin{tabularx}{\textwidth}{@{}l c c L@{}}
\toprule
\textbf{Checkpoint} & \textbf{Fast tokenizer} & \textbf{Offset mapping} & \textbf{Kết luận} \\
\midrule
\code{vinai/phobert-base} & \no & \no & Không dùng được cho trích xuất theo offset \\
\code{vinai/phobert-base-v2} & \no & \no & Không dùng được cho trích xuất theo offset \\
\code{nguyenvulebinh/vi-mrc-base} & --- & --- & Bị giới hạn truy cập (gated) trên HuggingFace \\
\code{FPTAI/videberta-base} & --- & --- & Không tìm thấy model id \\
\code{uitnlp/visobert} & \yes & \yes & \textbf{Được chọn} \\
\bottomrule
\end{tabularx}
\end{table}

Thiếu fast tokenizer là yếu tố chặn chứ không chỉ gây bất tiện: không có offset mapping thì
cách duy nhất có sẵn để lấy lại chuỗi đáp án là \code{decode()}, vốn không bảo toàn dấu
tiếng Việt và khoảng trắng gốc (Mục~\ref{sec:offsets}). PhoBERT còn yêu cầu đầu vào đã tách
từ, làm việc ánh xạ ngược về đoạn văn thô càng phức tạp. Về nguyên tắc có thể tự xây bảng
ánh xạ ký tự cho PhoBERT; nhóm đánh giá phần việc này vượt phạm vi thời gian của đồ án và
ghi nhận nó ở hướng phát triển. ViSoBERT được chọn thay thế vì là encoder tiếng Việt của
NLP@UIT --- nhóm công bố ViQuAD --- và có fast tokenizer. Hạn chế đã biết \emph{trước khi
chạy}: ViSoBERT được tiền huấn luyện trên văn bản mạng xã hội, khác miền Wikipedia của
ViQuAD.

\subsection{Gán nhãn phục vụ phân tích}
\label{sec:tagging}

\paragraph{Độ dài đoạn văn.} Mỗi câu hỏi được xếp vào một trong bốn nhóm theo số \emph{từ}
(tách theo khoảng trắng) của đoạn văn: \code{<100}, \code{100-200}, \code{200-300},
\code{300+}; biên trên của mỗi nhóm là đóng (200 từ thuộc nhóm \code{100-200}).

\paragraph{Phạm vi suy luận của câu hỏi.} UIT-ViQuAD \emph{không} cung cấp nhãn
single-sentence/multi-sentence, nên nhóm tự suy ra bằng heuristic: tách đoạn văn thành câu;
nếu tồn tại một câu đơn lẻ chứa ít nhất 60\% số token (âm tiết, không trùng lặp) của câu
hỏi thì câu hỏi được gán \code{single-sentence}, ngược lại \code{multi-sentence}. Chỉ câu hỏi
có đáp án được gán nhãn này, vì câu không có đáp án không có ``phạm vi bằng chứng'' để phân
loại.

\begin{warnbox}
\textbf{Giới hạn của heuristic.} Nhãn \code{question\_type} là xấp xỉ, không phải nhãn vàng.
Heuristic dựa trên \emph{độ trùng từ vựng}: một câu hỏi được diễn đạt lại bằng từ đồng nghĩa
có thể bị gán \code{multi-sentence} dù bằng chứng nằm trong một câu. Vì các mô hình --- đặc
biệt là TF-IDF --- cũng hưởng lợi từ độ trùng từ vựng, chênh lệch giữa hai nhóm phản ánh
\emph{một phần} độ khó do ít trùng từ, không chỉ độ khó của suy luận nhiều câu. Ngưỡng 0,6
được chọn theo quan sát và là một siêu tham số của heuristic.
\end{warnbox}

\subsection{Harness đánh giá và truy vết}
\label{sec:harness}

Hàm \code{run\_evaluation()} thực hiện theo thứ tự: kiểm tra \code{assert\_gradeable}; chạy bộ
dự đoán trên từng câu và đo độ trễ thực bằng \code{time.perf\_counter}; tính metric; gán nhãn
và tách nhóm; ghi kết quả. Mỗi tệp kết quả có cấu trúc như sau (minh hoạ giản lược của
\code{results/eval\_<mô hình>\_validation.json}; giá trị số là \emph{ví dụ}, số thật nằm
trong tệp JSON và trong các bảng ở Mục~\ref{sec:experiments}):

\begin{lstlisting}
{
  "model": "<tên mô hình>", "dataset": "UIT-ViQuAD 2.0",
  "split": "validation", "n": <n>,
  "timestamp": "<ISO 8601>", "commit": "<sha>", "device": "mps",
  "env": {"torch": "2.14.0", "platform": "macOS-26.4.1-arm64-arm-64bit"},
  "overall":         {"EM": <EM>, "F1": <F1>, "count": <n>},
  "answerable_only": {"EM": <EM>, "F1": <F1>, "count": <n_hasans>},
  "impossible_only": {"EM": <EM>, "count": <n_noans>},
  "avg_latency_ms": <ms>,
  "by_context_length": {"300+": {"EM": <EM>, "F1": <F1>, "count": <n>,
                                 "unreliable": true}, ...},
  "by_question_type":  {..., "_note": "Tổng của bảng này khác n vì câu không có đáp án bị loại khỏi phân tích loại câu hỏi"},
  "sample_predictions": [ ... ]
}
\end{lstlisting}

Hai cơ chế trong cấu trúc này nhắm vào lỗi trình bày: nhóm có $n < 30$ tự động mang cờ
\code{unreliable: true}, và khi tổng của một bảng tách nhóm không bằng $n$ (vì câu không có
đáp án bị loại khỏi phân tích loại câu hỏi), harness sinh trường \code{\_note} giải thích
lý do.

\subsection{Đăng ký giả thuyết trước khi chạy}
\label{sec:prereg}

Trước khi chạy bất kỳ đánh giá nào, nhóm ghi lại dự đoán định lượng và các ``tín hiệu báo
động'' vào \code{results/hypotheses.md} (ngày 12/9/2026), với cam kết không sửa sau khi thấy
kết quả. Mục đích là phòng vệ chống tự lừa mình: nếu không biết trước mình mong đợi gì, mọi
con số đều trông hợp lý --- kể cả con số sai. Script \code{check\_hypotheses.py} đối chiếu tự
động kết quả với các khoảng dự đoán lưu trong \code{results/hypotheses.json} và với các tín
hiệu báo động (Bảng~\ref{tab:alarms}); kết quả đối chiếu được trình bày ở
Mục~\ref{sec:hypothesis-check}.

\begin{table}[htbp]
\centering
\caption{Các tín hiệu báo động được định nghĩa để điều tra thay vì ăn mừng}
\label{tab:alarms}
\renewcommand{\arraystretch}{1.18}
\small
\begin{tabularx}{\textwidth}{@{}L L@{}}
\toprule
\textbf{Quan sát} & \textbf{Nghi vấn cần điều tra} \\
\midrule
EM của TF-IDF $> 10\%$ & Lỗi metric hoặc rò rỉ dữ liệu \\
EM của bất kỳ mô hình nào $> 85\%$ & Rò rỉ dữ liệu \\
Mô hình fine-tune kém hơn zero-shot & Lỗi cấu hình huấn luyện: lệch nhãn, tốc độ học, fp16 trên MPS \\
EM $\approx$ F1 (chênh $\le 0{,}5$) & Mô hình suy sụp về ``luôn trả rỗng'' \\
EM $\approx$ tỉ lệ câu không có đáp án (chênh $\le 1{,}0$) & Mô hình chỉ ghi điểm nhờ đoán rỗng \\
\bottomrule
\end{tabularx}
\end{table}

\noindent Ba tín hiệu đầu được ghi trong \code{hypotheses.md} trước khi chạy. Tệp này cũng ghi
trước một tín hiệu định tính: \emph{``điểm trên câu impossible $\approx$ 100\% mà answerable
$\approx$ 0\% $\Rightarrow$ mô hình học luôn trả về rỗng''}. Hai tín hiệu cuối của bảng là dạng
\emph{định lượng hoá} của tín hiệu đó, được bổ sung vào \code{hypotheses.json} sau khi quan
sát đường cong huấn luyện của ViSoBERT; các khoảng dự đoán định lượng không bị thay đổi.

\subsection{Kỹ thuật phần mềm và kiểm thử}
\label{sec:testing}

Mã nguồn được viết theo TDD qua 9 giai đoạn, bắt đầu từ metric và tầng dữ liệu, sau đó
đến baseline, cửa sổ trượt, harness, fine-tuning, hình vẽ và demo. Test được đánh dấu
\code{slow} khi cần tải tokenizer hoặc mô hình, để vòng lặp phát triển nhanh
(\code{pytest -m "not slow"}) chạy trong khoảng 2 giây. Bộ test được kiểm chứng ngược bằng
\emph{mutation testing}: cố tình phá mã nguồn và xác nhận có test thất bại, qua đó bịt các
lỗ hổng kiểm thử (Bảng~\ref{tab:modules}).

\begin{table}[htbp]
\centering
\caption{Module mã nguồn, số test và độ phủ (commit \code{c5b06ba})}
\label{tab:modules}
\renewcommand{\arraystretch}{1.12}
\small
\begin{tabularx}{\textwidth}{@{}l L r r@{}}
\toprule
\textbf{Module} & \textbf{Trách nhiệm} & \textbf{Số test (chậm)} & \textbf{Độ phủ$^\ast$} \\
\midrule
\code{normalize.py} & Chuẩn hoá chuỗi (quyết định A/B/C) & 15 & 100\% \\
\code{metrics.py} & EM, F1, tổng hợp theo quy ước SQuAD~2.0 & 25 & 100\% \\
\code{data.py} & Parse, chống rò rỉ, lấy mẫu con & 36 & 94\% \\
\code{tagging.py} & Nhóm độ dài, loại câu hỏi & 12 & 97\% \\
\code{baseline\_tfidf.py} & Baseline TF-IDF & 14 & 95\% \\
\code{windowing.py} & Cửa sổ trượt, chọn span, xác suất & 43 (10) & 76\% \\
\code{transformer\_qa.py} & Suy luận, bằng chứng cho quyết định & 26 (26) & 17\% \\
\code{features.py} & Gán nhãn token cho huấn luyện & 15 (15) & 0\% \\
\code{training.py} & Lịch học, chọn epoch, chẩn đoán & 34 & 94\% \\
\code{evaluate.py} & Harness, breakdown, truy vết & 16 & 92\% \\
\code{demo/} (6 module) & Logic ứng dụng web, tách khỏi Streamlit & 176 & --- \\
\code{reporting/} & Sinh hình và bảng cho báo cáo & 29 & --- \\
CLI / đóng gói & \code{test\_evaluation\_cli}, \code{test\_docker}, \code{test\_scripts} & 16 / 10 / 7 & --- \\
\midrule
\multicolumn{2}{@{}l}{\textbf{Tổng}: 423 test nhanh đều đạt; 51 test chậm cần tải mô hình} & \textbf{474 (51)} & \textbf{75\%} \\
\bottomrule
\end{tabularx}
\par\vspace{3pt}
{\footnotesize $^\ast$ Độ phủ đo khi chỉ chạy 423 test nhanh, phạm vi \code{src/mrc}. Module có test thuộc nhóm \code{slow} vì vậy có độ phủ \emph{đo được} thấp hơn phạm vi kiểm thử thực tế --- \code{transformer\_qa.py} có 26 test nhưng cả 26 đều cần tải mô hình thật. Hai thư mục \code{demo} và \code{reporting} nằm ngoài phạm vi đo nên không có số ở cột này.}
\end{table}

\subsection{Ứng dụng web}
\label{sec:demo-design}

Ứng dụng Streamlit \code{app/streamlit\_app.py} chỉ là lớp giao diện: logic nằm trong
\code{demo/logic.py} (hàm \code{answer()} và \code{highlight()}), nhận bộ dự đoán qua tham số nên
test được bằng đối tượng giả mà không cần chạy server hay tải mô hình. Người dùng chọn một
trong ba mô hình Transformer (chỉ mô hình đã có trên máy mới hiện ra), nhập đoạn văn và câu
hỏi, điều chỉnh ngưỡng $\tau$ trong khoảng $[-10, 10]$. Kết quả hiển thị đáp án, tô sáng vị
trí đáp án trong đoạn văn và độ trễ suy luận; khi mô hình kết luận không có đáp án, ứng dụng
giải thích rằng đây là hành vi hợp lệ với UIT-ViQuAD~2.0. Ảnh chụp giao diện được trình bày
ở Mục~\ref{sec:demo}.
